% from ega4-17-fr.tex
\oldpage[IV-4]{100}

que l'hypothèse entraîne que l'homomorphisme $(f^*(\Omega^1_{Y/k}))_x\to(\Omega^1_{X/k})_x$ est surjectif, donc, en remplaçant $X$ par un voisinage ouvert de $x$, on peut supposer que l'homomorphisme $f^*(\Omega^1_{Y/k})\to\Omega^1_{X/k}$ est surjectif (Bourbaki, \emph{Alg. comm.}, chap. II, § 5, n\textsuperscript{o} 1, prop. 2); par suite $f$ est \emph{non ramifié} (17.2.2). Nous allons voir d'abord qu'on peut se borner au cas où $x$ est \emph{rationnel sur $k$}. En effet, si l'on pose $k'=k(x)$, et $X'=X\otimes_k k'$, $Y'=Y\otimes_k k'=\operatorname{Spec}(k'[T_1,\ldots,T_n])$, il existe un point $x'\in X'$ au-dessus de $x$, tel que $k(x')=k'$. Pour prouver que $f$ est étale au point $x$, il suffit de montrer que $f'=f_{(k')}:X'\to Y'$ est étale au point $x'$ (17.7.1, (ii)); d'ailleurs $f'$ est non ramifié (17.3.3, (iii)) et l'on a $\dim_{x'}(X')=n$ (4.2.7). De la même manière on peut, en remplaçant $k'$ par une extension algébriquement close de $k'$, supposer que $k$ est \emph{algébriquement clos}. Posons alors $y=f(x)$, $A=\sh{O}_{Y,y}$, $B=\sh{O}_{X,x}$; comme le corps résiduel de $B$ est égal à $k$, il en est de même de celui de $A$, donc $x$ (resp. $y$) est un point fermé de $X$ (resp. $Y$) (I, 6.4.2) et l'on a par suite $\dim(A)=\dim(B)=n$. Puisque $f$ est non ramifié, l'homomorphisme $\widehat A\to\widehat B$ est surjectif (17.4.4, f'')); mais comme $\dim(\widehat A)=\dim(\widehat B)=n$, et que $\widehat A$, étant un anneau régulier, est intègre, l'homomorphisme $\widehat A\to\widehat B$ est aussi injectif (0, 16.3.10). Cet homomorphisme est donc bijectif, ce qui achève de prouver que $f$ est étale au point $x$ (17.6.3, e'')).

\begin{corollary}[17.15.4]
\label{IV.17.15.4}
Sous les hypothèses générales de (17.15.3), supposons en outre que $k(x)$ soit une extension finie et séparable de $k$ et que les germes $(g_i)_x$ appartiennent à $\mathfrak m_x$. Alors les conditions a), b), c) de (17.15.3) équivalent aussi à chacune des suivantes:
\begin{enumerate}
  \item[d)] Les $n$ germes $(g_i)_x$ engendrent l'idéal maximal $\mathfrak m_x$ de $\sh{O}_{X,x}$.
  \item[d')] L'anneau $\sh{O}_{X,x}$ est régulier et les $(g_i)_x$ forment un système régulier de paramètres pour cet anneau (0, 17.1.6).
\end{enumerate}
\end{corollary}

En effet, d') entraîne trivialement d). D'autre part, comme $k(x)$ est une extension finie et séparable de $k$, on a $\Omega^1_{k(x)/k}=0$ (0, 20.6.20) et $k(x)=\sh{O}_{X,x}/\mathfrak m_x$ est une $k$-algèbre formellement lisse pour les topologies discrètes (0, 19.6.1), donc la suite exacte (0, 20.5.14.1) s'applique à $A=k$, $B=\sh{O}_{X,x}$, $\mathfrak{R}=\mathfrak m_x$, et fournit un isomorphisme canonique
\[
\delta:\mathfrak m_x/\mathfrak m_x^2\xrightarrow{\sim}(\Omega^1_{X/k})_x\otimes_{\sh{O}_{X,x}}k(x).
\]
De là on déduit d'abord l'équivalence des conditions c) et d), compte tenu du lemme de Nakayama. D'autre part, si $f$ est étale au point $x$, l'anneau $\sh{O}_{X,x}$ est régulier et de dimension $n$, puisque $x$ est un point fermé de $X$ (I, 6.4.2), et $n$ éléments de $\mathfrak m_x$ qui engendrent cet idéal maximal forment alors nécessairement un système régulier de paramètres pour $\sh{O}_{X,x}$ (0, 17.1.6), ce qui prouve que a) entraîne d').

\begin{proposition}[17.15.5]
\label{IV.17.15.5}
Soient $k$ un corps, $X$ un préschéma localement de type fini sur $k$, $x$ un point de $X$, $n=\dim_x(X)$. Les conditions suivantes sont équivalentes:
\begin{enumerate}
  \item[a)] $X$ est lisse sur $k$ au point $x$.
  \item[b)] $X$ est différentiellement lisse sur $k$ au point $x$.
  \item[c)] $(\Omega^1_{X/k})_x$ est un $\sh{O}_{X,x}$-module libre de rang $n$.
  \item[d)] $(\Omega^1_{X/k})_x$ est un $\sh{O}_{X,x}$-module admettant un système de $n$ générateurs.

